\subsection*{5.3 Constructions of $\Xi,\alpha,\beta$}
\begin{proof}[Proof of Lemma 26]
Assume that $x_1,\ldots,x_k$ is a 0-instance $\mathrm{UDISJ}_{k,2^{\ell-k}-1}$. Let $i_0\in[2^{\ell-k}-1]$ be such that $x_{1,i_0}=x_{2,i_0}=\cdots=x_{k,i_0}=1$. Hereinafter $\gamma$ denotes a bijection from $[2^{\ell-k}-1]\cup\{*\}$ to $\{0,1\}^{\ell-k}$, $E$ denotes an invertible matrix from $\mathbb F_2^{\ell\times\ell}$ and $e_1,e_2,\ldots,e_\ell$ denote rows of $E$. Before presenting constructions of $\Xi,\alpha$, and $\beta$ we explain how we are going to establish symmetry and difference properties from Definition24. For every $s\in\{0,1\}^{\ell-k}$ and $b\in\{0,1\}$ we introduce the following notation:
\[
R(s,b,E):=\{(s,b,z)E\mid z\in\mathbb F_2^{k-1}\}.
\]
\begin{authorclaim}{Claim 27}
The last $2^{k+1}$ rows of $C_{\vec x}(\gamma,E)$ form $R(\gamma(*),0,E)\cup R(\gamma(*),1,E)$, each occurring twice. This is $\operatorname{Fake}(\gamma,E)$.
The rows with indices $[2^k(i-1)+1;2^ki]$ for $i\in[2^{\ell-k-1}]\setminus\{i_0\}$ are $R(\gamma(i),0,E)\cup R(\gamma(i),1,E)$, each occurring once.
The rows with indices $[2^k(i_0-1)+1;2^ki_0]$ are $R(\gamma(i_0),0,E)$, each occurring twice.
\end{authorclaim}
\begin{authorclaim}{Claim 28}
$R(s,b,E)=(\sum_{j=1}^{\ell-k}s_je_j+be_{\ell-k+1})+\operatorname{Span}(e_{\ell-k+2},\ldots,e_\ell)$.
\end{authorclaim}
\begin{proof}
\begin{align*}
R(s,b,E)&=\{(s,b,z)(e_1,\ldots,e_\ell)^T\mid z\in\mathbb F_2^{k-1}\}\\
&=\left\{\sum_{j=1}^{\ell-k}s_je_j+be_{\ell-k+1}+\sum_{i=1}^{k-1}z_ie_{\ell-k+1+i}\mid z\in\mathbb F_2^{k-1}\right\}\\
&=\left(\sum_{j=1}^{\ell-k}s_je_j+be_{\ell-k+1}\right)+\operatorname{Span}(e_{\ell-k+2},\ldots,e_\ell).
\end{align*}
\end{proof}
\begin{authorclaim}{Claim 29}
$|R(s,b,E)|=2^{k-1}$.
\end{authorclaim}
\begin{proof}
By Claim28, $|R(s,b,E)|=|(\sum_{j=1}^{\ell-k}s_je_j+be_{\ell-k+1})+\operatorname{Span}(e_{\ell-k+2},\ldots,e_\ell)|=|\operatorname{Span}(e_{\ell-k+2},\ldots,e_\ell)|=2^{k-1}$.
\end{proof}
\begin{authorclaim}{Claim 30}
The sets $R(s,b,E)$ for $s\in\{0,1\}^{\ell-k},b\in\{0,1\}$ are disjoint.
\end{authorclaim}
\begin{proof}
Consider two vectors $u\in R(s,b,E)$ and $v\in R(s',b',E)$ such that $(s,b)\ne(s',b')$. Then, by Claim28, $u$ and $v$ have different coordinates in the basis $e_1,e_2,\ldots,e_\ell$, hence $u\ne v$.
\end{proof}
\begin{authorclaim}{Claim 31}
For bijections $\gamma,\gamma'$ and invertible $E,E'$, if
\begin{align*}
&R(\gamma(i_0),0,E)\cup R(\gamma(*),0,E)\cup R(\gamma(*),1,E)\\
&\qquad=R(\gamma'(i_0),0,E')\cup R(\gamma'(*),0,E')\cup R(\gamma'(*),1,E'),\\
&R(\gamma(i_0),1,E)=R(\gamma'(i_0),1,E'),
\end{align*}
then $C_{\vec x}(\gamma,E)$ and $C_{\vec x}(\gamma',E')$ differ only by a row permutation.
\end{authorclaim}
\begin{proof}
By Claim27, rows from $R(\gamma(i_0),1,E)$ do not appear in $C_{\vec x}(\gamma,E)$, rows from $R(\gamma(i_0),0,E)\cup R(\gamma(*),0,E)\cup R(\gamma(*),1,E)$ appear in $C_{\vec x}(\gamma,E)$ exactly twice. The matrix $C_{\vec x}(\gamma,E)$ has $2^\ell+2^k$ rows. All rows of $C_{\vec x}(\gamma,E)$ that are not in $R(\gamma(i_0),1,E)\cup R(\gamma(*),0,E)\cup R(\gamma(*),1,E)$, by Claim27, appear in $C_{\vec x}(\gamma,E)$ exactly once.
By Claims29 and30, $|R(\gamma(i_0),0,E)\cup R(\gamma(*),0,E)\cup R(\gamma(*),1,E)|=3\cdot2^{k-1}$, hence, the number of rows of $C_{\vec x}(\gamma,E)$ that are not in $R(\gamma(i_0),1,E)\cup R(\gamma(*),0,E)\cup R(\gamma(*),1,E)$ equals $2^\ell-2^{k+1}$. By Claims29 and30, the number of $\ell$-bit strings not from $R(\gamma(i_0),1,E)\cup R(\gamma(i_0),0,E)\cup R(\gamma(*),0,E)\cup R(\gamma(*),1,E)$ is also $2^\ell-2^{k+1}$. Hence, all rows from $\{0,1\}^\ell\setminus(R(\gamma(i_0),0,E)\cup R(\gamma(*),0,E)\cup(\gamma(*),1,E)\cup R(\gamma(i_0),1,E))$ appear in $C_{\vec x}(\gamma,E)$ exactly once. Thus, matrices $C_{\vec x}(\gamma,E)$ and $C_{\vec x}(\gamma',E')$ have the same set of rows and each row appears the same number of times in each of these matrices.
\end{proof}
Write $\alpha(\gamma,E)=(\gamma_\alpha,E_\alpha)$ and $\beta(\gamma,E)=(\gamma_\beta,E_\beta)$. We construct them so that
\begin{equation}
\begin{aligned}
R(\gamma(i_0),1,E)&=R(\gamma_\alpha(i_0),1,E_\alpha)=R(\gamma_\beta(i_0),1,E_\beta),\\
R(\gamma(i_0),0,E)&=R(\gamma_\alpha(*),0,E_\alpha)=R(\gamma_\beta(*),0,E_\beta),\\
R(\gamma(*),1,E)&=R(\gamma_\alpha(*),1,E_\alpha)=R(\gamma_\beta(i_0),0,E_\beta),\\
R(\gamma(*),0,E)&=R(\gamma_\alpha(i_0),0,E_\alpha)=R(\gamma_\beta(*),1,E_\beta).
\end{aligned}\tag{2}
\end{equation}
Notice that by Claim31, equations(2) imply the symmetry property. Equations(2) also imply the difference property. Indeed, the three Fake sets are, respectively,
\begin{align*}
& R(\gamma(*),1,E)\cup R(\gamma(*),0,E),\\
& R(\gamma(*),1,E)\cup R(\gamma(i_0),0,E),\\
& R(\gamma(*),0,E)\cup R(\gamma(i_0),0,E),
\end{align*}
Hence, by Claim30, $\operatorname{Fake}(\gamma,E)\cap\operatorname{Fake}(\gamma_\alpha,E_\alpha)\cap\operatorname{Fake}(\gamma_\beta,E_\beta)=\varnothing$.
It remains to construct $\Xi,\alpha,\beta$ satisfying (2) and $\Pr[(\boldsymbol\gamma,\mathbf E)\in\Xi]>2/3+1/100$.
\paragraph{Definition of $\Xi$}
$(\gamma,E)\in\Xi$ iff there exist $m,n\in[\ell-k]$ with $\gamma(*)_m=1,\gamma(i_0)_m=0$ and $\gamma(*)_n=0,\gamma(i_0)_n=1$, i.e. $\gamma(*)$ and $\gamma(i_0)$ are coordinatewise incomparable.
$\boldsymbol\gamma(i_0)$ and $\boldsymbol\gamma(*)$ are uniform distinct elements. For independent uniform random subsets $\mathbf S,\mathbf T\subseteq[\ell-k]$,
\begin{align*}
\Pr[(\boldsymbol\gamma,\mathbf E)\in\Xi]
&=1-\Pr[\boldsymbol\gamma(i_0)\le\boldsymbol\gamma(*)\vee\boldsymbol\gamma(*)\le\boldsymbol\gamma(i_0)]\\
&\ge1-2\Pr[\boldsymbol\gamma(i_0)\le\boldsymbol\gamma(*)]\\
&=1-2\Pr[\mathbf S\subseteq\mathbf T\mid\mathbf S\ne\mathbf T]\\
&\ge1-2\Pr[\mathbf S\subseteq\mathbf T]\\
&=1-2\prod_{j=1}^{\ell-k}(1-\Pr[j\in\mathbf S\wedge j\notin\mathbf T])\\
&=1-2(3/4)^{\ell-k}>2/3+1/100\qquad(\ell-k\ge7).
\end{align*}
Hence largeness holds.
\paragraph{Construction of $\alpha$}
Let $E_\alpha$ have rows
\[
(e'_1,\ldots,e'_\ell)=
(e_1,\ldots,e_{\ell-k},e_{\ell-k+1}+\sum_{j=1}^{\ell-k}(\gamma(i_0)_j+\gamma(*)_j)e_j,e_{\ell-k+2},\ldots,e_\ell),
\]
and $\gamma_\alpha(i)=\gamma(*)$ for $i=i_0$, $\gamma(i_0)$ for $i=*$, and $\gamma(i)$ otherwise.
\begin{authorclaim}{Claim 32}
$\alpha:\Xi\to\Xi$ is bijective.
\end{authorclaim}
\begin{proof}
Notice that rows of $E'$ form a basis since $\sum_{j=1}^{\ell-k}(\gamma(i_0)_j+\gamma(*)_j)e_j\in\operatorname{Span}(e_1,\ldots,e_{\ell-k})$. The mapping $\gamma\mapsto\gamma_\alpha$ is bijective since it just swaps $\gamma(i_0)$ and $\gamma(*)$. Since the condition on $\gamma(i_0)$ and $\gamma(*)$ does not change after application of $\alpha$, we get that $\alpha(\Xi)\subseteq\Xi$. Notice that
\[
\sum_{j=1}^{\ell-k}(\gamma(i_0)_j+\gamma(*)_j)e_j
=\sum_{j=1}^{\ell-k}(\gamma_\alpha(i_0)_j+\gamma_\alpha(*)_j)e'_j,
\]
hence $\alpha(\gamma_\alpha,E_\alpha)=(\gamma,E)$, hence $\alpha$ is bijective.
\end{proof}
\begin{authorclaim}{Claim 33}
For all $(\gamma,E)\in\Xi$,
\begin{align*}
R(\gamma_\alpha(i_0),1,E_\alpha)&=R(\gamma(i_0),1,E),\\
R(\gamma_\alpha(i_0),0,E_\alpha)&=R(\gamma(*),0,E),\\
R(\gamma_\alpha(*),0,E_\alpha)&=R(\gamma(i_0),0,E),\\
R(\gamma_\alpha(*),1,E_\alpha)&=R(\gamma(*),1,E).
\end{align*}
\end{authorclaim}
\begin{proof}
Use Claim28 and $S:=\operatorname{Span}(e_{\ell-k+2},\ldots,e_\ell)=\operatorname{Span}(e'_{\ell-k+2},\ldots,e'_\ell)$:
\begin{align*}
R(\gamma_\alpha(i_0),1,E_\alpha)
&=(\sum_j\gamma_\alpha(i_0)_je'_j+e'_{\ell-k+1})+S\\
&=(\sum_j\gamma(*)_je'_j+e'_{\ell-k+1})+S
=(\sum_j\gamma(i_0)_je_j+e_{\ell-k+1})+S
=R(\gamma(i_0),1,E),\\
R(\gamma_\alpha(i_0),0,E_\alpha)
&=(\sum_j\gamma_\alpha(i_0)_je'_j)+S=(\sum_j\gamma(*)_je_j)+S=R(\gamma(*),0,E),\\
R(\gamma_\alpha(*),0,E_\alpha)
&=(\sum_j\gamma_\alpha(*)_je'_j)+S=(\sum_j\gamma(i_0)_je_j)+S=R(\gamma(i_0),0,E),\\
R(\gamma_\alpha(*),1,E_\alpha)
&=(\sum_j\gamma_\alpha(*)_je'_j+e'_{\ell-k+1})+S\\
&=(\sum_j\gamma(i_0)_je'_j+e'_{\ell-k+1})+S
=(\sum_j\gamma(*)_je_j+e_{\ell-k+1})+S=R(\gamma(*),1,E).
\end{align*}
Here each sum is over $j=1,\ldots,\ell-k$.
\end{proof}
\paragraph{Construction of $\beta$}
Set $\gamma_\beta=\gamma_\alpha$. Define $j_{\min}=\min\{j\in[\ell-k]:\gamma(*)_j=1\wedge\gamma(i_0)_j=0\}$, which exists on $\Xi$. Set $E_\beta=(e''_1,\ldots,e''_\ell)$ with
\[
e''_j=\begin{cases}
e_j,&j\notin\{j_{\min},\ell-k+1\},\\
\sum_{i=1}^{\ell-k}(\gamma(*)_i+\gamma(i_0)_i)e_i,&j=\ell-k+1,\\
e_{j_{\min}}+e_{\ell-k+1},&j=j_{\min}.
\end{cases}
\]
\begin{authorclaim}{Claim 34}
$\beta:\Xi\to\Xi$ is bijective.
\end{authorclaim}
\begin{proof}
Let us verify that $\beta$ is injective. Given $\gamma_\beta$ we can easily recover $\gamma$, hence we can recover $j_{\min}$ as well. Then
\begin{align*}
\sum_{i=1}^{\ell-k}(\gamma(i_0)_i+\gamma(*)_i)e''_i+e''_{\ell-k+1}
&=\sum_{i\in[\ell-k]\setminus\{j_{\min}\}}(\gamma(i_0)_i+\gamma(*)_i)e_i
+e_{j_{\min}}+e_{\ell-k+1}+e''_{\ell-k+1}\\
&=e_{\ell-k+1}+
\underbrace{\sum_{i\in[\ell-k]\setminus\{j_{\min}\}}(\gamma(i_0)_i+\gamma(*)_i)e_i+e_{j_{\min}}}_{e''_{\ell-k+1}}+e''_{\ell-k+1}\\
&=e_{\ell-k+1}.
\end{align*}
Thus, we can uniquely recover $e_{\ell-k+1}$ and, hence, also recover $e_{j_{\min}}=e''_{j_{\min}}+e_{\ell-k+1}$; for $j\in[\ell]\setminus\{j_{\min},\ell-k+1\}$, $e_j=e''_j$. Hence, $\beta$ is injective. Notice that since we represent $e_1,\ldots,e_\ell$ as linear combinations of $e''_1,\ldots,e''_\ell$, then $e''_1,\ldots,e''_\ell$ is a basis, hence the matrix $E_\beta$ is invertible. Thus, we verify that $\beta(\Xi)\subseteq\Xi$ and $\beta$ is injective, hence $\beta$ is bijective.
\end{proof}
\begin{authorclaim}{Claim 35}
For all $(\gamma,E)\in\Xi$,
\begin{align*}
R(\gamma_\beta(i_0),1,E_\beta)&=R(\gamma(i_0),1,E),\\
R(\gamma_\beta(i_0),0,E_\beta)&=R(\gamma(*),1,E),\\
R(\gamma_\beta(*),0,E_\beta)&=R(\gamma(i_0),0,E),\\
R(\gamma_\beta(*),1,E_\beta)&=R(\gamma(*),0,E).
\end{align*}
\end{authorclaim}
\begin{proof}
Put $S=\operatorname{Span}(e_{\ell-k+2},\ldots,e_\ell)=\operatorname{Span}(e''_{\ell-k+2},\ldots,e''_\ell)$. Recall $\gamma(*)_{j_{\min}}=1$, $\gamma(i_0)_{j_{\min}}=0$. Sums below run from1 to $\ell-k$:
\begin{align*}
R(\gamma_\beta(i_0),1,E_\beta)
&=\sum_i\gamma_\beta(i_0)_ie''_i+e''_{\ell-k+1}+S\\
&=\sum_i\gamma(*)_ie_i+e_{\ell-k+1}+e''_{\ell-k+1}+S\\
&=\sum_i\gamma(*)_ie_i+e_{\ell-k+1}+\sum_i(\gamma(*)_i+\gamma(i_0)_i)e_i+S\\
&=\sum_i\gamma(i_0)_ie_i+e_{\ell-k+1}+S=R(\gamma(i_0),1,E),\\
R(\gamma_\beta(i_0),0,E_\beta)
&=\sum_i\gamma_\beta(i_0)_ie''_i+S
=\sum_i\gamma(*)_ie_i+e_{\ell-k+1}+S=R(\gamma(*),1,E),\\
R(\gamma_\beta(*),0,E_\beta)
&=\sum_i\gamma_\beta(*)_ie''_i+S=\sum_i\gamma(i_0)_ie_i+S=R(\gamma(i_0),0,E),\\
R(\gamma_\beta(*),1,E_\beta)
&=\sum_i\gamma_\beta(*)_ie''_i+e''_{\ell-k+1}+S
=\sum_i\gamma(i_0)_ie_i+e''_{\ell-k+1}+S
=\sum_i\gamma(*)_ie_i+S=R(\gamma(*),0,E).
\end{align*}
\end{proof}
Claims33 and35 imply equations(2).
\end{proof}
